% Primera lectura: del plano a una superficie, esfera y elipsoide, sistemas, marcos y grillas
\section{Pasar del plano a una superficie}\label{cap:06}\label{sec:plano-superficie}

{\small\vspace{1mm}\linkUnidad{10_jacobiano_primera_forma}{Simulación interactiva: Unidad 14 — Jacobiano y primera forma fundamental}\par}


Sobre una superficie curva no existe una única base cartesiana válida en todas partes. En
cada punto aparece un plano tangente y una base local que cambia al desplazarse. La métrica
indica cómo convertir pequeños cambios de coordenadas en distancias y ángulos reales sobre
la superficie.

\paragraph{Primera traducción algebraica.}
\[
  ds^2=E\,du^2+2F\,du\,dv+G\,dv^2.
\]
Los coeficientes \(E,F,G\) describen cuánto miden y qué ángulo forman los dos vectores de la
base local. La expresión transforma pequeños cambios \((du,dv)\) en una distancia geométrica.

\begin{idea}[De lo global a lo local]
Una superficie curva se estudia mediante pequeños parches que se parecen a un plano. En cada
parche se usan coordenadas locales; al conectar los parches se recupera la geometría global.
Esta es la intuición básica de una variedad.
\end{idea}

\begin{idea}[En palabras del profesor]
``La derivada es la mejor aproximación lineal a una función en determinado punto''; ``el
gradiente es perpendicular a una superficie/curva de nivel [\ldots] porque el gradiente es
una forma lineal, y la curva de nivel 0 de esa forma lineal es el plano tangente'' (resumen de clase 9, cursada 2025).
\end{idea}

\section{Esfera, elipsoide y representación plana}\label{sec:esfera-primera}

{\small\vspace{1mm}\linkUnidad{12_curvatura}{Simulación interactiva: Unidad 16 — Curvatura como forma cuadrática}\par}
{\small\vspace{1mm}\linkUnidad{13_elipsoide}{Simulación interactiva: Unidad 17 — Geometría del elipsoide}\par}


La Tierra se aproxima primero mediante una esfera y luego mediante un elipsoide de revolución.
Su curvatura cambia con la posición, por lo que una diferencia angular no representa la misma
distancia en todas las latitudes.

Una proyección cartográfica lleva esa superficie al plano. Como una superficie curva no puede
desarrollarse perfectamente sobre un plano, toda proyección debe elegir qué propiedad
privilegiar: ángulos, áreas o determinadas distancias.

\begin{idea}[En palabras del profesor]
``Buscar una intuición de rotaciones y coordenadas sobre la cáscara de la esfera, que va
a terminar conectando con el sistema de coordenadas de latitud y longitud'' (resumen de clase 2, cursada 2025).
\end{idea}

\section{Sistemas, marcos y grillas}\label{sec:sistemas-marcos}

Un sistema de coordenadas define reglas matemáticas para expresar posiciones. Un marco de
referencia materializa esas reglas en una época concreta mediante puntos, estaciones y
observaciones.

Cuando se cambia de marco pueden bastar traslaciones, rotaciones y escala. Si existen
deformaciones regionales, se utilizan grillas de corrección que asignan desplazamientos
variables según la ubicación. En datos ráster, el cambio de coordenadas obliga además a
remuestrear porque los centros de píxel no coinciden exactamente.

\begin{resultado}[Mapa mental de la materia]
Las primeras transformaciones enseñan a reconocer invariantes. Las afinidades amplían la
familia permitiendo deformaciones. La geometría diferencial lleva esas ideas a superficies
curvas, y la geodesia las aplica al elipsoide, las proyecciones y los marcos de referencia.
\end{resultado}

\begin{idea}[En palabras del profesor]
``Una operación muy importante es la composición [\ldots]: en general es necesario para
construir un sistema de coordenadas que haya conmutatividad a la hora de componer
movimientos por coordenadas'' (resumen de clase 2, cursada 2025).
\end{idea}

